\documentclass[11pt]{article}

\usepackage{fontspec}
\setmainfont{Noto Sans CJK KR}
\usepackage{amsmath,amssymb,mathtools}
\usepackage[a4paper,margin=16mm,top=15mm,bottom=16mm]{geometry}
\usepackage{xcolor}
\usepackage[most]{tcolorbox}
\usepackage{enumitem}
\usepackage{titlesec}

\setlist[itemize]{leftmargin=1.4em,itemsep=0.25em,topsep=0.35em}
\setlist[enumerate]{leftmargin=1.5em,itemsep=0.3em,topsep=0.35em}
\setlength{\parskip}{0.35em}
\setlength{\parindent}{0pt}
\allowdisplaybreaks

\titleformat{\section}{\large\bfseries\color{blue!45!black}}{\thesection.}{0.5em}{}
\titleformat{\subsection}{\normalsize\bfseries\color{blue!35!black}}{\thesubsection}{0.45em}{}
\titlespacing*{\section}{0pt}{1.15em}{0.45em}
\titlespacing*{\subsection}{0pt}{0.85em}{0.3em}

\tcbset{breakable,enhanced,sharp corners,boxsep=2.5pt,left=5pt,right=5pt,top=3pt,bottom=3pt}

\newtcolorbox{mainbox}[1]{
  colback=blue!6,colframe=blue!55!black,title={주요\quad #1},
  fonttitle=\bfseries,coltitle=white,attach boxed title to top left={yshift=-2mm,xshift=2mm},
  boxed title style={colback=blue!55!black,sharp corners}
}
\newtcolorbox{confbox}[1]{
  colback=orange!8,colframe=orange!75!black,title={혼동\quad #1},
  fonttitle=\bfseries,coltitle=white,attach boxed title to top left={yshift=-2mm,xshift=2mm},
  boxed title style={colback=orange!75!black,sharp corners}
}
\newtcolorbox{relbox}[1]{
  colback=violet!7,colframe=violet!60!black,title={관계\quad #1},
  fonttitle=\bfseries,coltitle=white,attach boxed title to top left={yshift=-2mm,xshift=2mm},
  boxed title style={colback=violet!60!black,sharp corners}
}
\newtcolorbox{fillbox}{
  colback=green!7,colframe=green!45!black,
  title={이어서 계산},
  fonttitle=\bfseries,coltitle=white,attach boxed title to top left={yshift=-2mm,xshift=2mm},
  boxed title style={colback=green!45!black,sharp corners}
}

\newcommand{\ex}[1]{\par\textbf{세부\quad #1}\par}
\renewcommand{\contentsname}{차례}

\begin{document}

\begin{center}
{\LARGE\bfseries 공업수학 단권화 노트}\\[0.35em]
{\large 1계 미분방정식 \textperiodcentered\ 정적분}\\[0.45em]
{\small 수업 필기를 주제별로 다시 묶은 정리. 쪽 순서가 아니라 같은 풀이끼리 모았다.}
\end{center}

\vspace{0.2em}
수업 노트에 적힌 공식과 예제만 사용했다. 낙서, 지운 메모, 과제만 적혀 있고 풀이가 없는 식은 빼었다. 식이 중간에 끊기거나 마지막 계산이 식과 맞지 않는 곳은 녹색 상자에서 같은 공식으로 끝까지 계산했다.

\tableofcontents
\vspace{0.4em}

\section{어떤 식으로 풀지}

\begin{mainbox}{식의 모양과 이름}
노트에 나온 1계 방정식은 아래 다섯 가지다. 이름을 고르기 전에, \(y'\)의 계수와 우변에 \(y\)가 어떻게 들어 있는지를 본다.
\[
\begin{array}{lll}
\hline
\text{식} & \text{노트에서의 이름} & \text{해의 꼴}\\
\hline
g(y)\,y'=f(x) & \text{변수분리} & \text{양쪽을 적분}\\[0.25em]
y'+a(x)y=0 & \text{동차 선형} & y=Ce^{-\int a(x)\,dx}\\[0.25em]
y'+a(x)y=f(x) & \text{비동차 선형} &
y=e^{-\int a}\,\!\Bigl[\int f\,e^{\int a}\,dx+C\Bigr]\\[0.35em]
y'+a(x)y=f(x)\,y^{n} & \text{베르누이} & u=y^{1-n}\text{으로 선형화}\\[0.25em]
M\,dx+N\,dy=0,\; \partial M/\partial y=\partial N/\partial x & \text{완전} & f(x,y)=C\\
\hline
\end{array}
\]
\end{mainbox}

\begin{relbox}{다섯 공식이 서로 어떻게 겹치는가}
\begin{itemize}
\item 동차 선형은 비동차 공식에서 \(f=0\)인 경우다. 적분 안이 상수만 남는다.
\item 베르누이 방정식은 그 자체로 \(y'+a(x)y=f(x)\) 꼴이 아니다. \(u=y^{1-n}\)으로 바꾸면 비동차 선형이 되고, 그때 위의 일반해를 쓴다.
\item 완전미분은 \(y'=\cdots\)로 정리하기 전의 꼴 \(M\,dx+N\,dy=0\)이다. \(\partial M/\partial y\neq\partial N/\partial x\)이면 완전이 아니므로, 그 식을 \(y'\)에 대해 풀어 선형인지 다시 본다.
\item 변수분리는 \(y\)쪽과 \(x\)쪽이 곱으로 갈라질 때 쓴다. 선형 공식의 \(a(x)\) 자리에 \(y\)가 들어가면 선형이 아니다.
\end{itemize}
\end{relbox}

노트 여백에만 단어로 적혀 있는 것: 미분, 상미분, 편미분, 계수, 차수, 고계 미분방정식, 선형 미분방정식, 비선형 미분방정식. 이 단어들의 정의는 노트에 따로 없다.

\section{변수분리}

\begin{mainbox}{변수분리}
\[
g(y)\,\frac{dy}{dx}=f(x)
\quad\Longrightarrow\quad
\int g(y)\,dy=\int f(x)\,dx+C.
\]
\(y\)만 있는 인수를 왼쪽으로, \(x\)만 있는 인수를 오른쪽으로 옮긴 뒤 적분한다.
\end{mainbox}

\ex{예제 1-1 (a)}
\[
\frac{dy}{dx}+\frac{y}{x}=0.
\]
\[
\frac{dy}{y}=-\frac{dx}{x},\qquad
\ln|y|=-\ln|x|+C_1,\qquad
\ln|xy|=C_1.
\]
\[
y=\frac{C}{x}.
\]
\(C\)는 \(0\)이 아닌 상수와 \(y=0\)을 포함한다. 같은 식을 뒤의 예제 1-6 (a)에서 동차 선형 공식으로 다시 푼다.

\ex{예제 1-1 (b)}
\[
\frac{dy}{dx}=\frac{2}{x^{2}-1}=\frac{2}{(x-1)(x+1)}.
\]
\[
\frac{2}{(x-1)(x+1)}=\frac{1}{x-1}-\frac{1}{x+1},
\]
\[
y=\ln|x-1|-\ln|x+1|+C=\ln\left|\frac{x-1}{x+1}\right|+C.
\]

\ex{초기조건이 있는 변수분리}
\[
\frac{dy}{dx}=\frac{y}{x+2},\qquad y(0)=1.
\]
\[
\ln|y|=\ln|x+2|+C_0.
\]
\(y(0)=1\)이면 \(0=\ln 2+C_0\)이므로 \(C_0=-\ln 2\).
\[
\ln|y|=\ln\left|\frac{x+2}{2}\right|.
\]
\(y(0)=1>0\)이고 \(x=0\) 근처에서 \(x+2>0\)이므로
\[
y=\frac{x+2}{2}.
\]

\[
\frac{dy}{dx}=e^{-x}-1,\qquad y(0)=1.
\]
\[
y=-e^{-x}-x+C.
\]
\(1=-1+C\)에서 \(C=2\).
\[
y=-e^{-x}-x+2.
\]

\ex{노트에 \(y^{2}\)이 들어 있어 분리한 식}
\[
\frac{dy}{dx}+(x+1)y^{2}=0.
\]
\[
y^{-2}\,dy=-(x+1)\,dx,\qquad
-\frac{1}{y}=-\frac{x^{2}}{2}-x+C.
\]
\[
y=\frac{1}{\dfrac{x^{2}}{2}+x+C}.
\]
\(n=2\)인 베르누이 방정식이기도 하다. 노트에서는 치환 없이 변수분리로 풀었다.

\ex{\(\dfrac{1}{x}\dfrac{dy}{dx}=y^{2}-y,\; y(1)=2\)}
\[
\frac{dy}{y(y-1)}=x\,dx.
\]
\[
\frac{1}{y(y-1)}=\frac{1}{y-1}-\frac{1}{y},
\]
\[
\ln\left|\frac{y-1}{y}\right|=\frac{x^{2}}{2}+C.
\]
\(x=1,\;y=2\)이면 \(\ln(1/2)=1/2+C\)이므로
\[
C=\ln\frac{1}{2}-\frac{1}{2}.
\]

\begin{fillbox}
노트는 상수 \(C\)까지 구하고 멈춘다. \(C\)를 도로 넣으면
\[
\ln\left|\frac{y-1}{y}\right|
=\frac{x^{2}}{2}+\ln\frac{1}{2}-\frac{1}{2}.
\]
\(y(1)=2\)이므로 \((y-1)/y>0\)인 쪽을 택한다.
\[
\frac{y-1}{y}=\frac12\exp\!\left(\frac{x^{2}-1}{2}\right),
\qquad
y=\left(1-\frac12\exp\!\left(\frac{x^{2}-1}{2}\right)\right)^{-1}.
\]
\(x=1\)에서 \(y=2\)가 된다.
\end{fillbox}

\begin{confbox}{\(1/(y(y-1))\)의 부분분수}
\[
\frac{1}{y(y-1)}=\frac{1}{y-1}-\frac{1}{y}
\]
이지, \(\dfrac{1}{y}+\dfrac{1}{y-1}\)이 아니다. 부호를 같게 두면 \(\ln|y(y-1)|\)이 되어, 초기조건을 넣은 줄의 \(\ln|(y-1)/y|\)와 다른 식이 된다.
\end{confbox}

\section{1계 선형}

\begin{mainbox}{정의 1-2}
\[
\frac{dy}{dx}+a(x)y=f(x).
\]
\begin{itemize}
\item \(f(x)=0\)이면 동차, \(f(x)\neq 0\)이면 비동차.
\item 노트에 적힌 보기: \(y'+x^{2}y=e^{x}\)는 비동차, \(y'+x^{2}y=0\)은 동차.
\end{itemize}
동차일 때
\[
\frac{dy}{y}=-a(x)\,dx,\qquad
\ln|y|=-\int a(x)\,dx+C_0,
\]
\[
y=Ce^{-\int a(x)\,dx}.
\]
비동차의 일반해는
\begin{equation}
y=e^{-\int a(x)\,dx}
\left[\int f(x)\,e^{\int a(x)\,dx}\,dx+C\right].
\label{eq:general}
\end{equation}
해 앞에 붙는 지수는 \(-\int a\), 적분 안의 지수는 \(+\int a\)이다.
\end{mainbox}

\begin{confbox}{동차와 비동차, 그리고 지수 부호}
\begin{itemize}
\item 동차·비동차는 우변 \(f\)가 \(0\)인지로 가른다. \(a(x)\)가 상수인지와는 별개다.
\item \(e^{-\int a}\)와 \(e^{\int a}\)를 같은 자리에 넣으면 적분이 풀리지 않거나 다른 해가 나온다.
\item \(f=e^{x}\)이면 적분 안이 \(1\)이 되는 경우가 있고, \(f=e^{2x}\)이면 적분 안이 \(e^{x}\)가 된다. 우변의 지수와 \(\int a\)의 부호를 같이 본다.
\end{itemize}
\end{confbox}

\ex{예제 1-6 (a), 동차}
\[
\frac{dy}{dx}+\frac{y}{x}=0,\qquad a(x)=\frac{1}{x},\qquad \int a(x)\,dx=\ln|x|.
\]
\[
y=Ce^{-\ln|x|}=\frac{C}{|x|}.
\]
예제 1-1 (a)와 같은 식이다. \(0\)을 포함하지 않는 구간에서는 \(y=C/x\)로 쓸 수 있다.

\ex{예제 1-6 (b), 동차, \(y(0)=1\)}
\[
\frac{dy}{dx}-\frac{1}{(x+1)(x+2)}y=0,
\qquad
a(x)=-\frac{1}{(x+1)(x+2)}.
\]
\[
\frac{1}{(x+1)(x+2)}=\frac{1}{x+1}-\frac{1}{x+2},
\]
\[
\int a(x)\,dx=\ln\left|\frac{x+2}{x+1}\right|,
\qquad
y=C\,\frac{x+1}{x+2}.
\]
\begin{fillbox}
노트는 일반해에서 멈춘다. \(y(0)=1\)이면 \(C/2=1\)이므로 \(C=2\).
\[
y=\frac{2(x+1)}{x+2}.
\]
\end{fillbox}

\ex{예제 1-7 (a), 비동차}
\[
\frac{dy}{dx}-y=e^{2x}.
\]
\(a(x)=-1\), \(\int a\,dx=-x\), \(f(x)=e^{2x}\). 식 \eqref{eq:general}에 넣으면
\[
y=e^{x}\left[\int e^{2x}\cdot e^{-x}\,dx+C\right]
=e^{x}\left[\int e^{x}\,dx+C\right]
=e^{2x}+Ce^{x}.
\]

\begin{confbox}{이 예의 적분}
노트 중간에는 적분 안이 \(1\)이 되고 \(y=e^{x}(x+C)\)까지 적혀 있다. \(y=e^{x}(x+C)\)는 \(y'-y=e^{x}\)의 해다. 우변이 \(e^{2x}\)이면 적분 안은 \(e^{2x}e^{-x}=e^{x}\)이고, 해는 \(e^{2x}+Ce^{x}\)이다.
\end{confbox}

\ex{예제 1-7 (b), 비동차, \(y(1)=1\)}
\[
\frac{dy}{dx}=\frac{y+x}{x}
\quad\Longleftrightarrow\quad
\frac{dy}{dx}-\frac{1}{x}y=1.
\]
\(a(x)=-1/x\), \(\int a\,dx=-\ln|x|\), \(f=1\).
\[
y=|x|\bigl(\ln|x|+C\bigr).
\]
\(x=1,\;y=1\)이면 \(C=1\). 이 초기조건이 있는 쪽 \(x>0\)에서는
\[
y=x(\ln x+1).
\]

\ex{동차, \(y(-1)=e^{-1}\)}
\[
\frac{dy}{dx}+\frac{1}{x^{2}}y=0,\qquad a(x)=\frac{1}{x^{2}}.
\]
\[
y=Ce^{-\int x^{-2}\,dx}=Ce^{1/x}.
\]
\(e^{-1}=Ce^{-1}\)이므로 \(C=1\).
\[
y=e^{1/x}.
\]

\ex{\(y'+y=\cos x\)}
\(a(x)=1\), \(f(x)=\cos x\).
\[
y=e^{-x}\left[\int\cos x\cdot e^{x}\,dx+C\right].
\]
노트는 부분적분을
\[
u=\cos x,\; dv=e^{x}\,dx,\qquad
du=-\sin x\,dx,\; v=e^{x}
\]
까지 두고, \(-\sin x\cdot e^{x}+\cos x\cdot e^{x}\)에서 끊긴다.

\begin{fillbox}
\(\displaystyle\int e^{x}\cos x\,dx\)를 두 번 부분적분한다.
\[
\int e^{x}\cos x\,dx
=e^{x}\cos x+\int e^{x}\sin x\,dx
=e^{x}\cos x+e^{x}\sin x-\int e^{x}\cos x\,dx.
\]
\[
\int e^{x}\cos x\,dx=\frac{e^{x}}{2}(\sin x+\cos x).
\]
\[
y=e^{-x}\left[\frac{e^{x}}{2}(\sin x+\cos x)+C\right]
=\frac12(\sin x+\cos x)+Ce^{-x}.
\]
검산: \(y'+y=\cos x\).
\end{fillbox}

\section{베르누이 방정식}

\begin{mainbox}{1.4 베르누이}
\[
\frac{dy}{dx}+a(x)y=f(x)\,y^{n}.
\]
\[
u=y^{1-n}
\]
으로 두면 \(u\)의 1계 선형 방정식이 된다. 그 해를 구한 뒤 \(y\)로 되돌린다.
\end{mainbox}

\begin{confbox}{치환의 지수}
\(u=y^{n}\)이 아니라 \(u=y^{1-n}\)이다. \(n=2\)이면 \(u=y^{-1}\), \(n=-1\)이면 \(u=y^{2}\)이다. \(u\)로 바꾼 식의 \(a(x)\)는 원래 \(a(x)\)와 부호·계수가 달라진다. 나눈 뒤에 다시 \(u'\)으로 정리한 결과를 선형 공식에 넣는다.
\end{confbox}

\ex{1-10}
\[
\frac{dy}{dx}+\frac{1}{x}y=xy^{2}.
\]
\(n=2\), \(u=y^{-1}\). \(y^{2}\)으로 나누고 \(u'=-y^{-2}y'\)를 쓰면
\[
u'-\frac{1}{x}u=-x.
\]
\(x>0\)에서 \(\int a\,dx=-\ln x\)이므로 식 \eqref{eq:general}은
\[
u=x\left[\int -1\,dx+C\right]=x(-x+C)=-x^{2}+Cx.
\]
\[
y=\frac{1}{Cx-x^{2}}.
\]

\ex{1-11}
\[
\frac{dy}{dx}-y=xy^{-1}.
\]
\(n=-1\), \(u=y^{2}\), \(u'=2yy'\).
\[
u'=2y(y+xy^{-1})=2y^{2}+2x=2u+2x,
\]
\[
u'-2u=2x.
\]
\(a=-2\), \(f=2x\).
\[
u=e^{2x}\left[\int 2xe^{-2x}\,dx+C\right].
\]
\[
\int 2xe^{-2x}\,dx=-e^{-2x}\Bigl(x+\frac12\Bigr),
\]
\[
u=-(x+\tfrac12)+Ce^{2x}.
\]
\[
y=\pm\sqrt{Ce^{2x}-x-\frac12}.
\]

\ex{\(y'+y=e^{-x}y^{2}\)}
노트에는 식과 ``베르누이 방정식''이라는 이름만 있다.

\begin{fillbox}
\(n=2\), \(u=y^{-1}\). \(y^{2}\)으로 나누면
\[
y^{-2}y'+y^{-1}=e^{-x}.
\]
\(y^{-2}y'=-u'\)이므로
\[
u'-u=-e^{-x}.
\]
\(a=-1\), \(f=-e^{-x}\).
\[
u=e^{x}\left[\int -e^{-x}\cdot e^{-x}\,dx+C\right]
=e^{x}\left[\frac12 e^{-2x}+C\right]
=\frac12 e^{-x}+Ce^{x}.
\]
\[
y=\frac{1}{Ce^{x}+\dfrac12 e^{-x}}
=\frac{2e^{x}}{2Ce^{2x}+1}.
\]
\end{fillbox}

\begin{relbox}{1-10과 바로 위의 식}
둘 다 \(n=2\)라서 \(u=y^{-1}\)이다. 차이는 원래의 \(a(x)\)와 \(f(x)\)뿐이다.
\[
\begin{array}{c|cc}
 & a(x) & f(x)\\
\hline
\text{1-10} & 1/x & x\\
y'+y=e^{-x}y^{2} & 1 & e^{-x}
\end{array}
\]
\(u\) 방정식의 계수는 이 값 그대로가 아니라, \(y^{2}\)으로 나눈 뒤 \(u'=-y^{-2}y'\)로 부호를 맞춘 결과다.
\end{relbox}

\section{완전미분방정식}

\begin{mainbox}{1.5 완전}
\[
M(x,y)\,dx+N(x,y)\,dy=0.
\]
\begin{enumerate}
\item \(\partial M/\partial y\)와 \(\partial N/\partial x\)를 계산한다. 같으면 완전이다.
\item \(f\)를 찾는다. 노트에서는 \(\partial f/\partial x=M\)로 두고 \(x\)로 적분한 뒤, \(\partial f/\partial y=N\)으로 \(h(y)\)를 정한다.
\item 해는 \(f(x,y)=C\).
\end{enumerate}
\(M\)은 \(dx\)의 계수, \(N\)은 \(dy\)의 계수다.
\end{mainbox}

\begin{confbox}{\(M\)과 \(N\)을 바꾸면 완전성 판정이 뒤집힌다}
\(\partial/\partial y\)는 \(M\)에, \(\partial/\partial x\)는 \(N\)에 작용한다. \(h\)는 먼저 적분하지 않은 변수의 함수다. \(x\)로 먼저 적분하면 \(h=h(y)\)이고, \(h'(y)\)는 \(\partial f/\partial y\)와 \(N\)을 비교해서 나온다.
\end{confbox}

\ex{1-12 (a)}
\[
(2xy+1)\,dx+x^{2}\,dy=0.
\]
\[
M=2xy+1,\; N=x^{2},\qquad
\frac{\partial M}{\partial y}=2x=\frac{\partial N}{\partial x}.
\]
\[
f=\int(2xy+1)\,dx=x^{2}y+x+h(y),
\]
\[
\frac{\partial f}{\partial y}=x^{2}+h'(y)=x^{2}
\quad\Rightarrow\quad h'(y)=0.
\]
\[
x^{2}y+x=C,\qquad y=\frac{C-x}{x^{2}}.
\]
두 식은 같은 해다. 뒤의 식은 지운 것이 아니다.

\ex{1-12 (b), \(x=0,\;y=1\)}
\[
y^{2}\,dx+2(xy+1)\,dy=0.
\]
\[
M=y^{2},\; N=2xy+2,\qquad
\frac{\partial M}{\partial y}=2y=\frac{\partial N}{\partial x}.
\]
\[
f=\int y^{2}\,dx=xy^{2}+h(y),
\]
\[
\frac{\partial f}{\partial y}=2xy+h'(y)=2xy+2
\quad\Rightarrow\quad h'(y)=2,\; h(y)=2y.
\]
\[
xy^{2}+2y=C.
\]
\(x=0,\;y=1\)이면 \(C=2\).
\[
xy^{2}+2y=2.
\]

\ex{\((2xy+1)\,dx+(x^{2}+x)\,dy=0\)}
노트에는 식과 ``완전미분방정식''이라는 이름만 있다.

\begin{fillbox}
노트의 완전성 조건으로 검사하면
\[
M=2xy+1,\; N=x^{2}+x,
\]
\[
\frac{\partial M}{\partial y}=2x,\qquad
\frac{\partial N}{\partial x}=2x+1.
\]
둘이 같지 않으므로 완전이 아니다. \(y'\)에 대해 풀면 1계 선형이 된다.
\[
\frac{dy}{dx}+\frac{2}{x+1}y=-\frac{1}{x(x+1)}.
\]
\(a(x)=2/(x+1)\), \(\int a\,dx=2\ln|x+1|\). \(x>-1\)에서
\[
\frac{d}{dx}\bigl[y(x+1)^{2}\bigr]
=-(x+1)^{2}\cdot\frac{1}{x(x+1)}
=-1-\frac{1}{x}.
\]
\[
y(x+1)^{2}=-x-\ln|x|+C,
\qquad
y=\frac{C-x-\ln|x|}{(x+1)^{2}}.
\]
\end{fillbox}

\section{정적분}

정적분 쪽은 치환, 우함수·기함수, 삼각함수 항등식, 부분분수, 부분적분이 노트에 있다. 하나의 적분에 두 가지가 같이 쓰이면, 먼저 쓰는 도구로 묶었다.

\subsection{치환}

\begin{mainbox}{정적분의 치환}
\[
\int f\bigl(g(x)\bigr)\,g'(x)\,dx=\int f(u)\,du.
\]
정적분에서는 \(x\)의 양 끝을 \(u\)의 양 끝으로 바꾼다.
\end{mainbox}

\ex{\(\displaystyle\int_{1}^{2}\frac{1}{(3-5x)^{2}}\,dx\)}
\(u=3-5x\), \(du=-5\,dx\). \(x=1\)이면 \(u=-2\), \(x=2\)이면 \(u=-7\).
\[
\int_{-2}^{-7}u^{-2}\Bigl(-\frac{du}{5}\Bigr)
=\frac15\int_{-7}^{-2}u^{-2}\,du
=\frac15\left[-\frac{1}{u}\right]_{-7}^{-2}
=\frac15\left(\frac12-\frac17\right)
=\frac{1}{14}.
\]

\ex{\(\displaystyle\int_{0}^{a}x\sqrt{x^{2}+a^{2}}\,dx\quad(a>0)\)}
\(u=x^{2}+a^{2}\), \(du=2x\,dx\). \(x=0\)이면 \(u=a^{2}\), \(x=a\)이면 \(u=2a^{2}\).
\[
\frac12\int_{a^{2}}^{2a^{2}}u^{1/2}\,du
=\frac13\Bigl[(2a^{2})^{3/2}-(a^{2})^{3/2}\Bigr]
=\frac{a^{3}}{3}(2\sqrt{2}-1).
\]

\ex{\(\displaystyle\int_{0}^{1/2}\frac{\sin^{-1}x}{\sqrt{1-x^{2}}}\,dx\)}
\(u=\sin^{-1}x\)이면 \(du=dx/\sqrt{1-x^{2}}\). \(x=0\)이면 \(u=0\), \(x=1/2\)이면 \(u=\pi/6\).
\[
\int_{0}^{\pi/6}u\,du=\frac12\left(\frac{\pi}{6}\right)^{2}=\frac{\pi^{2}}{72}.
\]

\subsection{우함수와 기함수}

\begin{mainbox}{대칭 구간}
\(f\)가 우함수, 즉 \(f(-x)=f(x)\)이면
\[
\int_{-a}^{a}f(x)\,dx=2\int_{0}^{a}f(x)\,dx.
\]
\(f\)가 기함수, 즉 \(f(-x)=-f(x)\)이면
\[
\int_{-a}^{a}f(x)\,dx=0.
\]
구간이 \(-a\)에서 \(a\)일 때만 이렇게 줄인다.
\end{mainbox}

\begin{confbox}{2배를 쓸지 0으로 볼지}
우함수 공식은 적분을 \(2\)배한다. 기함수 공식은 \(0\)이다. 노트에는 \(\tan x/(1+x^{2}+x^{4})\)를 \(2\int_{0}^{1}\)로 쓰기 시작했다가, 분자는 기함수, 분모의 \(1\), \(x^{2}\), \(x^{4}\)는 각각 우함수라고 고쳐 적혀 있다. 기함수를 우함수 공식에 넣지 않는다.
\end{confbox}

\ex{\(\displaystyle\int_{-1}^{1}\frac{\tan x}{1+x^{2}+x^{4}}\,dx\)}
\(\tan(-x)=-\tan x\)이고 분모는 \(x\)를 \(-x\)로 바꿔도 같다. 따라서 피적분함수는 기함수다.
\[
\int_{-1}^{1}\frac{\tan x}{1+x^{2}+x^{4}}\,dx=0.
\]
\([-1,1]\)에서 분모는 \(0\)이 아니고, \(\tan x\)의 불연속점 \(\pi/2\)도 이 구간 밖에 있다.

\ex{\(\displaystyle\int_{-\pi/4}^{\pi/4}(x^{3}+x^{4}\tan x)\,dx\)}
\begin{itemize}
\item \(x^{3}\)은 기함수.
\item \(x^{4}\)은 우함수, \(\tan x\)는 기함수이므로 \(x^{4}\tan x\)는 기함수.
\item 기함수끼리 더하면 기함수.
\end{itemize}
구간이 \(-\pi/4\)에서 \(\pi/4\)로 대칭이므로 적분은 \(0\)이다. \(\pi/4<\pi/2\)라서 \(\tan x\)는 이 구간에서 연속이다.

\subsection{삼각함수 적분}

\begin{mainbox}{노트에 별표로 적힌 항등식}
\[
\begin{aligned}
\cos^{2}x+\sin^{2}x&=1,\\
1+\tan^{2}x&=\sec^{2}x,\\
1+\cot^{2}x&=\csc^{2}x,\\
\cos^{2}x&=\frac{1+\cos 2x}{2},\\
\sin^{2}x&=\frac{1-\cos 2x}{2}.
\end{aligned}
\]
\(\cot x=\cos x/\sin x\). \(\cot^{2}x\) 적분에서는 \(\cot^{2}x=\csc^{2}x-1\)을 쓴다.
\end{mainbox}

\ex{\(\displaystyle\int_{0}^{\pi/2}\sin^{2}x\cos^{2}x\,dx\)}
노트는 반각으로 전개하다가
\[
\frac18\int_{0}^{\pi/2}(2-\cos 4x)\,dx
\]
에서 값이 없이 끝난다.

\begin{fillbox}
같은 반각 공식으로 정리하면
\begin{align*}
\sin^{2}x\cos^{2}x
&=\frac{1-\cos 2x}{2}\cdot\frac{1+\cos 2x}{2}
=\frac{1-\cos^{2}2x}{4}\\
&=\frac{\sin^{2}2x}{4}
=\frac{1-\cos 4x}{8}.
\end{align*}
\(1-\cos^{2}2x=(1-\cos 4x)/2\)이므로, 마지막 분자에 남는 것은 \(2-\cos 4x\)가 아니라 \(1-\cos 4x\)다.
\[
\int_{0}^{\pi/2}\sin^{2}x\cos^{2}x\,dx
=\frac18\int_{0}^{\pi/2}(1-\cos 4x)\,dx
=\frac18\cdot\frac{\pi}{2}
=\frac{\pi}{16}.
\]
\end{fillbox}

\ex{\(\displaystyle\int_{\pi/6}^{\pi/3}\cot^{2}x\,dx\)}
\[
\cot^{2}x=\csc^{2}x-1,
\]
\[
\int_{\pi/6}^{\pi/3}\cot^{2}x\,dx
=\bigl[-\cot x-x\bigr]_{\pi/6}^{\pi/3}.
\]
\[
\Bigl(-\cot\frac{\pi}{3}-\frac{\pi}{3}\Bigr)
-\Bigl(-\cot\frac{\pi}{6}-\frac{\pi}{6}\Bigr)
=\Bigl(-\frac{1}{\sqrt{3}}-\frac{\pi}{3}\Bigr)
-\Bigl(-\sqrt{3}-\frac{\pi}{6}\Bigr).
\]

\begin{fillbox}
\[
\sqrt{3}-\frac{1}{\sqrt{3}}-\frac{\pi}{6}
=\frac{3-1}{\sqrt{3}}-\frac{\pi}{6}
=\frac{2\sqrt{3}}{3}-\frac{\pi}{6}.
\]
노트 마지막 줄은 \(\sqrt{3}-\pi/3\)이다. 그 줄에는 \(-\cot(\pi/3)=-1/\sqrt{3}\)과 구간의 차 \(-\pi/3-(-\pi/6)=-\pi/6\)이 빠져 있다.
\end{fillbox}

\subsection{부분분수}

\ex{\(\displaystyle\int_{0}^{1}\frac{2}{2x^{2}+3x+1}\,dx\)}
\[
2x^{2}+3x+1=(2x+1)(x+1),
\]
\[
\frac{2}{(2x+1)(x+1)}=\frac{4}{2x+1}-\frac{2}{x+1}.
\]
\[
\int\left(\frac{4}{2x+1}-\frac{2}{x+1}\right)dx
=2\ln|2x+1|-2\ln|x+1|.
\]
\[
\bigl[2\ln(2x+1)-2\ln(x+1)\bigr]_{0}^{1}
=2\ln 3-2\ln 2
=2\ln\frac{3}{2}.
\]

\subsection{부분적분}

\begin{mainbox}{정적분의 부분적분}
노트에 적힌 식은
\[
\int_{a}^{b}f(x)g'(x)\,dx
=\Bigl[f(x)g(x)\Bigr]_{a}^{b}
-\int_{a}^{b}g(x)f'(x)\,dx.
\]
\end{mainbox}

\ex{\(\displaystyle\int_{0}^{1}\tan^{-1}x\,dx\)}
노트에는 적분만 있다.

\begin{fillbox}
\(f=\tan^{-1}x\), \(g'=1\)로 두면 \(f'=1/(1+x^{2})\), \(g=x\).
\begin{align*}
\int_{0}^{1}\tan^{-1}x\,dx
&=\bigl[x\tan^{-1}x\bigr]_{0}^{1}
-\int_{0}^{1}\frac{x}{1+x^{2}}\,dx\\
&=\frac{\pi}{4}-\frac12\bigl[\ln(1+x^{2})\bigr]_{0}^{1}
=\frac{\pi}{4}-\frac12\ln 2.
\end{align*}
\end{fillbox}

\ex{\(\displaystyle\int_{\sqrt{\pi/2}}^{\sqrt{\pi}}\theta^{3}\cos(\theta^{2})\,d\theta\)}
노트에는 적분만 있다.

\begin{fillbox}
\(u=\theta^{2}\)이면 \(du=2\theta\,d\theta\), \(\theta^{3}\,d\theta=\theta^{2}\cdot\theta\,d\theta=u\,du/2\).
\(\theta=\sqrt{\pi/2}\)이면 \(u=\pi/2\), \(\theta=\sqrt{\pi}\)이면 \(u=\pi\).
\[
\frac12\int_{\pi/2}^{\pi}u\cos u\,du.
\]
\(u\)와 \(\cos u\)를 부분적분하면 \(\int u\cos u\,du=u\sin u+\cos u\).
\begin{align*}
\frac12\bigl[u\sin u+\cos u\bigr]_{\pi/2}^{\pi}
&=\frac12\Bigl[(\pi\cdot 0+\cos\pi)
-\bigl(\tfrac{\pi}{2}\cdot 1+\cos\tfrac{\pi}{2}\bigr)\Bigr]\\
&=\frac12\Bigl[-1-\frac{\pi}{2}\Bigr]
=-\frac12-\frac{\pi}{4}.
\end{align*}
\end{fillbox}

\subsection{넓이}

\ex{타원 \(\dfrac{x^{2}}{a^{2}}+\dfrac{y^{2}}{b^{2}}=1\)로 둘러싸인 영역의 넓이}
노트에는 문제만 있다.

\begin{fillbox}
제1사분면의 \(y\)는 \(y=b\sqrt{1-x^{2}/a^{2}}\)이고, 대칭이므로 넓이는 이 곡선의 4배다.
\[
S=4b\int_{0}^{a}\sqrt{1-\frac{x^{2}}{a^{2}}}\,dx.
\]
\(x=a\sin\theta\)이면 \(dx=a\cos\theta\,d\theta\). \(x=0\)일 때 \(\theta=0\), \(x=a\)일 때 \(\theta=\pi/2\).
\[
S=4ab\int_{0}^{\pi/2}\cos^{2}\theta\,d\theta
=4ab\int_{0}^{\pi/2}\frac{1+\cos 2\theta}{2}\,d\theta
=2ab\Bigl[\theta+\frac12\sin 2\theta\Bigr]_{0}^{\pi/2}
=\pi ab.
\]
\(\cos^{2}\theta\)는 노트에 적힌 반각 공식 \((1+\cos 2\theta)/2\)다.
\end{fillbox}

\begin{relbox}{정적분 도구가 한 문제에서 겹칠 때}
\begin{itemize}
\item \(\theta^{3}\cos(\theta^{2})\)는 먼저 치환으로 \(u\cos u\)를 만들고, 그다음 부분적분한다.
\item 타원 넓이는 먼저 \(x=a\sin\theta\)로 치환하고, 남은 \(\cos^{2}\theta\)에 반각 공식을 쓴다.
\item 대칭 구간의 기함수 적분은 피적분함수를 전개하기 전에 \(0\)인지 확인한다. \(x^{3}+x^{4}\tan x\)가 그 경우다.
\end{itemize}
\end{relbox}

\end{document}
